\title{X-LoRA: Mixture of Low-Rank Adapter Experts, a Flexible Framework for Large Language Models with Applications in Protein Mechanics and Molecular Design}

\begin{abstract}
We present a mixture of expert methodology to develop fine-tuned large language models by employing a deep layer-wise token-level approach grounded in low-rank adaptation (LoRA). Beginning with a collection of pre-trained LoRA adapters, our gating mechanism leverages hidden states to dynamically blend adapted layers, enabling the X-LoRA model to tap into diverse functionalities and generate novel deep layer-wise combinations for task resolution. This design draws inspiration from biological concepts of universality and diversity, where neural network components are repurposed across various hierarchical configurations. Consequently, the X-LoRA model can be seamlessly integrated with any existing large language model (LLM) without requiring changes to the fundamental architecture. We create a specialized X-LoRA model offering scientific functionalities such as forward and inverse analysis tasks along with improved reasoning abilities, particularly targeting biomaterial analysis, protein mechanics, and design. This work provides access to models that are easily extensible and adaptable, possessing robust domain expertise and the ability to synthesize knowledge across disciplines. Incorporating experts in biology, mathematics, reasoning, bio-inspired materials, mechanics and materials science, chemistry, protein biophysics, mechanics, and quantum mechanics-based molecular properties, we perform a series of physics-centered case studies. These include assessments of knowledge recall, protein mechanics forward/inverse problems, protein design, adversarial agentic modeling with ontological knowledge graph construction, and molecular design. The model not only makes quantitative predictions of nanomechanical protein properties or quantum mechanical molecular characteristics but also interprets the outcomes and accurately identifies probable mechanisms underlying distinct molecular behaviors.
\end{abstract}

\section{Introduction}  
Large language models (LLMs)~\cite{Vaswani2017AttentionNeedc,Touvron2023LlamaModels,OpenAI2023GPT-4Report,Chowdhery2022PaLM:Pathways,Jiang2023Mistral7B,Gunasekar2023TextbooksNeed} have become highly popular, including their use in creating specialized models that excel in specific types of tasks, reasoning processes, or scientific areas~\cite{Bubeck2023SparksGPT-4,Buehler2023MechGPTModalities,Nejjar2023LLMsAnalysis,Buehler2023GenerativeDesign,Luu2023BioinspiredLLM:Materials,Luu2023GenerativeSolvents,Buehler2023MeLMProblemsb,Ge2023OpenAGI:Experts}. However, the training of these models can be expensive, particularly when a broad range of capabilities is required. Approaches such as low-rank adapters (LoRA)~\cite{hu2021lora} have been introduced as more resource-efficient alternatives, though these adaptations typically concentrate on more specialized knowledge domains. The core idea behind LoRA is to incorporate low-rank matrices added to the original full-sized matrices, designating these low-rank matrices as the sole trainable components. Because only the adapter layers undergo training, these models generally retain the information learned during the base model’s pre-training while tailoring the model to specific tasks, all while maintaining computational efficiency. Given the efficiency of training LoRA adapters, it becomes feasible to develop a large collection of specialized models through this method.  

Despite the efficiency of LoRA adapters in enabling specialized capabilities, a key challenge remains in how to combine multiple such adapters to create enhanced models that provide a cohesive, and potentially superior, set of functionalities. Indeed, the idea of mixing LLMs has garnered interest as a research topic~\cite{Kim2023SOLARUp-Scaling}. Experimental investigations, such as those involving spherically interpolated model parameters (SLERP) and related techniques~\cite{Arcee-ai/mergekit:Models.}, have demonstrated promising outcomes and suggest the potential of model combination to yield novel capabilities. Models formed through these mixing methods are usually computed {\it a priori} based on specific algorithms, and while they have exhibited encouraging results on certain benchmarks, further exploration is required to systematically develop mixed models and understand how distinct features arise.

We propose that dynamically mixing parameters could serve as a more versatile technique, as it enables the mixed model to investigate novel combinations during inference.

Similar methods have employed mixture of experts (MoE) frameworks~\cite{Jacobs1991AdaptiveExperts,Hampshire1992TheRecognition,Jordan1994HierarchicalAlgorithm}, including the recently introduced Mixtral LLM~\cite{Jiang2024MixtralExperts}. Nonetheless, these methods can be computationally intensive and require considerable resources even during inference. In contrast, we introduce a computationally efficient mixture-of-expert approach that utilizes a collection of distinct LoRA adapters, which can be straightforwardly trained independently. Our method embodies a high-granularity and adaptable framework, drawing inspiration from a biological model that repurposes neural network building blocks to create a hierarchy of interacting models, spanning from the foundational model to agentic systems (Fig.~\ref{fig:hierarchical_concept}). To evaluate its effectiveness quantitatively, our study concentrates on applying this approach to science-oriented LLMs, particularly focusing on protein mechanics challenges within the wider context of bio-inspired materials analysis and design.

We present an algorithm that dynamically combines adapters at the token-level state with fine granularity, permitting mixing at the level of individual layers. This method, termed X-LoRA, grants access to a variety of capabilities essential for scientific applications involving multimodal language modeling~\cite{Hu2023DeepScience,Buehler2023AMetamaterials,Buehler2022ProteinNetworks,}.

\subsection{Fundamental concepts of X-LoRA}

We begin with a brief overview of the principles underlying the development of LoRA adapters. The fundamental idea behind low-rank adaptation (LoRA)~\cite{hu2021lora} posits that weight updates possess a low “intrinsic dimension” and leverages this by keeping the original weights $W_0 \in \mathbb{R}^{d \times k}$ fixed while restricting updates to a low-rank decomposition

\begin{equation}
W_0 + \Delta W_0 = W_0 + BA 
\end{equation}

where $B \in \mathbb{R}^{d \times r}$ and $A \in \mathbb{R}^{r \times d}$ with the rank $r \ll {\rm min}(d,k)$. In the training process with LoRA~\cite{hu2021lora}, the original pretrained weights $W_0$ remain fixed, while only the matrices $A$ and $B$ are updated as trainable parameters. The forward pass for a LoRA layer can be represented by the equation:

\begin{equation}
    h = W_0x + \Delta W_0x = W_0x + BAx
\end{equation}

It should be noted that, in practical applications, LoRA adapters incorporate a fixed scalar factor $\alpha$. This scalar is applied to the decomposition matrices, resulting in

\begin{equation}
    h = W_0x + BAx \cdot \alpha
\end{equation}

Extending this concept, X-LoRA introduces a set of LoRA adapters, each endowed with distinct capabilities, as demonstrated in prior research focused on scientific LLMs (e.g., \cite{Buehler2023GenerativeDesign}). Rather than statically mixing or combining models, we employ a dynamic gating mechanism that adjusts the scaling of each LoRA adapter on a per-token and per-layer basis, enabling integration deeper within the model. Specifically, for a given LoRA adapter $i$, the original parameter $\alpha_i$ is modified by a scaling factor $\lambda_i$ to produce a new scaling parameter $\alpha^*_i$:

\begin{equation}
    \alpha^*_i = \alpha_i \cdot \lambda_i
\end{equation}

The scaling factor $\lambda_i$ is generated by an X-LoRA scaling head that leverages the model’s hidden states, serving as the trainable element of the X-LoRA framework. Because the scaling head can utilize rich contextual encodings, it can be efficiently implemented as a simple feed-forward neural network composed of one or a few layers (in X-LoRA’s implementation, users have the flexibility to define the structure of this network).

This scaling operation is performed individually for each token in the input sequence, thereby offering fine-grained control as each LoRA adapter active at specific layers within the LLM is scaled accordingly.

From this point onward, we describe this method of gating the adapter components through such multiplicative adjustments as scaling. Additional details regarding this technique are provided in the Materials and Methods section.

\subsection{Paper outline}

The structure of this paper is outlined as follows. Initially, we describe the methodology and training strategy, including the relevant datasets, which lead to the creation of an X-LoRA model endowed with specialized capabilities in the physical sciences, with a particular emphasis on biomaterials such as proteins. Subsequently, we present a sequence of experiments where we utilize the X-LoRA model across various tasks including question answering, conversational and agent-based modeling, protein design and analysis, among others. We perform an in-depth examination of the scaling behaviors and validate our approach by comparing it with molecular modeling and other physical datasets and methodologies.

\section{Results and discussion}

The development process for the X-LoRA model involves the following stages:

\begin{enumerate}
    \item Train a foundational base LLM (completed in prior work) 
    \item Train a series of individual adapters to build expertise across several domains (these represent distinct or overlapping skill sets and knowledge areas) 
    \item Train the combined X-LoRA model
\end{enumerate}

Our experimental work begins with the training of several LoRA adapters. We create nine adapters, each fine-tuned to specialize in specific expertise, based on the Zephyr-7B-$\beta$ model~\cite{HuggingFaceH4/zephyr-7b-betaFace}, which itself was developed atop the Mistral-7B model\cite{Jiang2023Mistral7B}:

\begin{enumerate}
    \item Bioinspired materials 
    \item Chain-of-thought (CoT) and reasoning 
    \item Chemistry 
    \item Mathematics 
    \item Physics 
    \item Biology 
    \item Mechanics and materials 
    \item Platypus: Logic and reasoning 
    \item Protein mechanics tasks (generative sequence-to-property and inverse capabilities)
\end{enumerate}

These focus areas and areas of expertise were selected because our ultimate goal is to develop a STEM-oriented, materials science-focused LLM. Nevertheless, we hypothesize that our method can be broadly applied to various fields, particularly for creating modalities that facilitate connections between diverse domains of knowledge~\cite{Buehler2023MechGPTModalities,Buehler2023MeLMProblemsb}.

Once the LoRA adapters have been generated, we proceed to train the X-LoRA model. For this purpose, we utilize a subset drawn from the combined pool of training samples originally used to train each individual LoRA adapter (for more information, see Materials and Methods, which includes detailed descriptions of the dataset, implementation, and related aspects).

We employ the X-LoRA framework as an experimental tool to evaluate its effectiveness across multiple scenarios. Our investigations encompass a range of tests, such as addressing complex questions, leveraging the generative model for protein design, conducting analysis and reasoning, and devising more advanced interaction strategies through adversarial agentic modeling. Throughout these assessments, we compare the performance of the base model against the X-LoRA model and monitor how the scaling mechanism engages different layers from various experts during answer generation or development.

\subsection{Question answering and observed X-LoRA layer-wise scaling weights}

Initially, we explore the distinctions between the X-LoRA model and the foundational base model. We pose the following question to both models:

\begin{quote}
\small\texttt{What material behavior (stiffening vs softening) causes supersonic crack speeds in mode II cracks? Think step by step.}
\end{quote}

The base model responds as follows:

\begin{quote}
\small\texttt{Supersonic crack speeds, also known as Rayleigh waves, occur in mode II cracks due to the stiffening material behavior of the material near the crack tip. In mode II fracture, the loading is shear-dominated and the crack propagates perpendicular to the direction of applied stress. Here's a step-by-step explanation:}

As the crack tip nears the material, the stresses around it escalate sharply. This creates a zone of intense stress concentration known as the stress singularity.

Within this stress singularity, the material undergoes elastic deformation that results in increased stiffness, referred to as stiffening behavior. This stiffening arises because the material adjacent to the crack tip becomes increasingly constrained as the crack advances [...]}
\end{quote}

Drawing from our domain expertise~\cite{Buehler2008AtomisticFailure}, this explanation is not only inaccurate but also unnecessarily verbose. Conversely, the X-LoRA model provides the following response:

\begin{quote}
\small\texttt{The material behavior that causes supersonic crack speeds in mode II cracks is stiffening.}
\end{quote}

This response is accurate, concise, and straightforward~\cite{Buehler2008AtomisticFailure}.

The contrast between the base model and the X-LoRA model is even more pronounced in the specialized tasks related to protein analysis and design. The task presented is:

\begin{quote}
\small\texttt{CalculateForceHistory< M S K S N T Y R M L V L L E D D T K I N K E D E K F L K G K P G K M H E F V D E L I L P F N V D E L D E L N T W F D K F D A E I C I P Y V E I L K E S G M K > }
\end{quote}

The reply from the base foundation model is:

\begin{quote}
\small\texttt{I'm sorry but the input provided is not a valid mathematical expression. It appears to be a sequence of characters that do not form a calculable function. Please provide a valid mathematical expression for me to calculate.

However, I can tell you that the sequence of characters you provided is actually a series of function and variable names that may be part of a larger program or script in a specific programming language. Without context, it's impossible to determine the intended calculation.

If you provide more information, such as the programming language and the context in which this expression appears, I may be able to provide more insight into its meaning and possible calculations. }
\end{quote}

The X-LoRA model responds as follows, successfully accomplishing the design task of producing the force-deformation curve corresponding to this protein:

\begin{quote}
\small\texttt{[0.013,0.066,0.078,0.045,0.072,0.126,0.127,0.107,0.119, 
0.074,0.106,0.098,0.103,0.128,
0.077, 0.069,0.131,0.108,0.144,0.174, 0.166,0.183,0.185,0.217,0.175,0.230,0.187,0.224, 0.267,0.301,0.315]}
\end{quote}

Figure~\ref{fig:QA_weights_sample_0} presents outcomes from this question-answering interaction alongside the observed X-LoRA scaling weights. The two queries—one centered on mechanics and materials (Fig.~\ref{fig:QA_weights_sample_0}a) and the other on protein mechanics (Fig.~\ref{fig:QA_weights_sample_0}b)—exhibit distinct activation patterns. Notably, we observe that a complex arrangement of scaling values is employed in each scenario, indicating that the X-LoRA model leverages a heterogeneous combination of different adapters across layers. Additionally, the heatmap in Fig.~\ref{fig:QA_weights_sample_0}a demonstrates a wider distribution of activation across numerous experts and layers, with certain regions showing strong activation. Conversely, the heatmap in Fig.~\ref{fig:QA_weights_sample_0}b reveals a more focused activation pattern, where high values are less diffusely spread. We interpret this as the gating function on the right being more selective in its expert activations.

In Fig.~\ref{fig:QA_weights_sample_0}, the bright line-like regions correspond to paths of strong activation, indicating that the scaling function selectively enhances certain experts for specific layers. The observed patterns imply that importance is not evenly distributed among all experts and layers; rather, it is generally sparse and selective. Such sparsity aligns with findings from other mixtures-of-experts models, where only a limited subset of experts is engaged for a given input to specialize in different aspects of the problem space.

By comparing these maps across various tasks, both the sparsity and the patterning may provide valuable insights into which experts the model depends on for different input types or at various processing stages within the neural network (we will later discuss how these maps evolve dynamically during the handling of complex tasks).

To deepen the analysis, we investigate a broader range of questions to understand how the X-LoRA model employs its different experts when presented with diverse prompts. Our primary interest is to determine whether, and if so, how, the X-LoRA model blends multiple adapters for questions that span overlapping skills or domains. Fig.~\ref{fig:QA_weights} presents the results from question answering alongside the observed X-LoRA scaling weights, with the queries detailed in Table~\ref{tab:table_qa}. In Table~\ref{tab:table_qa}, we include both the X-LoRA responses and those from the base model, Zephyr-7B-$\beta$. The distinctions between the two models are pronounced.

To highlight this clearly, we indicate incorrect or questionable responses using \redhighlight{red highlight}. This comparison reveals that X-LoRA achieves substantially higher accuracy and produces more succinct answers.

As expected, we observe intricate mixing of adapters and frequently the activation of several leading LoRA experts. For example, Fig.~\ref{fig:QA_weights}d shows that a question about pine cone mechanics triggers strong engagement of the mechanics/materials adapter along with the CoT/reasoning and bioinspired adapters. Although X-LoRA selects the correct letter, the answer is not fully accurate since pine cones open in dry rather than wet conditions. A subsequent question probing whether the release occurs in dry or wet conditions yields the correct response.

For a protein-related question, Fig.~\ref{fig:QA_weights}f demonstrates adapter mixing in reply to an inquiry about Bombyx mori silk mechanics, activating a combination of the bioinspired, mechanics/materials, and biology adapters. In clearly specialized tasks such as protein design (Fig.~\ref{fig:QA_weights}a), we observe predominantly singular activation of one expert.

A more detailed examination of the heatmap of scalings shown in Fig.~\ref{fig:QA_weights}a reveals that beyond the clearly significant role of the protein mechanics expert, there are multiple bands where particular experts receive considerable importance across various layers. This is especially evident for the mechanics/materials and reasoning experts. These bands are not continuous and exhibit a fluctuating pattern, which may indicate that these experts effectively handle specific conditions or features. Notably, the reasoning expert (located immediately to the left of the protein mechanics expert) stands out in several layers, particularly in the higher layers (approximately 100 to 140), where it achieves the highest activation. We suggest that this may point to an important contribution to the model’s final decision-making or output, although further investigation is needed to confirm this. Additionally, there is a more localized but pronounced activation for the physics and biology experts in certain layers, potentially reflecting a specialization of these experts for features or representations found in those regions.

Moreover, Fig.~\ref{fig:QA_token_history} illustrates a token-level history based on the summed scalings across all layers. This analysis uncovers histories for two examples: Fig.~\ref{fig:QA_token_history}a relates to a task involving pine cone seed release, while Fig.~\ref{fig:QA_token_history}b corresponds to a sequence of protein mechanics tasks. The protein mechanics example consisted of a series of prompt-answer exchanges: beginning with two property calculations and concluding with a generative task to design a new protein. Although expert utilization remains generally consistent, noticeable variations occur in the engagement of the LoRA experts throughout the process. Specifically, the protein mechanics expert becomes highly prominent in the final generation stages when the new protein is being designed.

Table~\ref{tab:table_qa} presents various use cases across different fields and compares performance against the base model alone. We provide a more detailed discussion of several of these cases below.

For example, we present a question combining protein science and biology with logical reasoning, focusing on protein expression patterns. X-LoRA provides the correct answer. In contrast, the base model's response contains some errors and logical flaws when evaluating the marker protein expressions in each cell type. The base model accurately states that fibroblasts express Protein B and do not express Protein A due to the mutual exclusivity of these two proteins. However, it fails to explicitly mention that fibroblasts also express Protein C. Since there is no rule prohibiting the co-expression of Protein C with Protein B, and fibroblasts have no restriction against Protein C, it is logical that they express Protein C as well. Regarding Neurons, the base model begins correctly by indicating that neurons express Protein A and not Protein B, based on the mutual exclusivity rule. Nevertheless, its analysis regarding the possibility of neurons expressing Protein C or both Protein C and Protein B is unclear and contains an error. Given that neurons express Protein A but not Protein B, the only remaining option consistent with the rules is that they express Protein C. There is no scenario allowing neurons to express Protein B, contrary to what the initial analysis suggested. Therefore, neurons express Protein A and Protein C, but not Protein B. For Epithelial cells, the base model correctly identifies that these cells do not express Protein A but do express both Protein B and Protein C. However, the explanation is overly complicated and includes incorrect claims about mutual exclusivity between Protein A and Protein C, which was not part of the original rules. The straightforward logic is that epithelial cells do not express Protein A but must express the other two proteins, Protein B and Protein C, since the initial instructions only restrict the co-expression of Protein A and Protein B.

X-LoRA, on the other hand, accurately determines the marker proteins expressed by each cell type by analyzing the given observations and logical inferences. Regarding fibroblasts, X-LoRA correctly infers that they do not express Protein A because the presence of Protein A excludes Protein B expression. Since fibroblasts do express Protein B, and considering that each cell type expresses a unique combination of three marker proteins (which was an incorrect assumption introduced, as the original statement did not specify the number of proteins each cell type expresses but only that they have a unique combination), the conclusion that fibroblasts also express Protein C is justified based on the provided information (apart from the misunderstanding about the protein count). For neurons, the model accurately identifies that they express Protein A and Protein C. The reasoning is valid, noting that Protein B cannot co-exist with Protein A, and because neurons express Protein A plus one other marker protein (not Protein B), they must express Protein C. Lastly, for epithelial cells, the X-LoRA model correctly concludes that they express Protein B and Protein C, since they do not express Protein A and, by elimination, as they express two marker proteins, these must be Protein B and Protein C.

In the inquiry regarding pine cone release under humid versus dry conditions, the base model provides a convoluted and factually incorrect response. Conversely, X-LoRA delivers a clear and well-reasoned answer. Similar examples include the question about the exceptional mechanical properties of Bombyx mori silk, where the base model incorrectly claims that silk fibers conduct electricity. The base model also struggles with several other domain-specific questions, such as the one concerning the role of hyperelasticity in maximum fracture speeds, where X-LoRA offers a concise and accurate response while the base model fails to provide an answer. In a question about the strong adhesion of gecko feet, X-LoRA accurately highlights the significance of setae and nano-scale effects, whereas the base model incorrectly attributes the process to special purpose glue proteins. These and other examples detailed in the table and throughout the paper provide compelling evidence of X-LoRA’s superior capabilities relative to the base model.

Figure~\ref{fig:perf_comp_bioinspiredexam} presents a comparison of knowledge recall evaluation tests using the bio-inspired knowledge test set introduced in \cite{Luu2023BioinspiredLLM:Materials}. It is evident that X-LoRA delivers the highest performance, despite being a considerably smaller model than both the BioinspiredLLM and Llama-BioLLM models (7B parameters for X-LoRA compared to 13B parameters in the others). Additional comparisons with recently published LLMs, specifically Orca-13B and Llama-13b-chat, are also included.

The figure further contrasts X-LoRA with the base Zephyr-7B-$\beta$ model, demonstrating a marked improvement in performance by X-LoRA. These comparisons suggest that strong results can be obtained using X-LoRA even when smaller base models are employed.

Figure~\ref{fig:image_synth} illustrates an application in image synthesis, where X-LoRA is used to generate a prompt intended for use with other models such as DALL-E 3, as exemplified here.

\begin{quote}
\small\texttt{Create a prompt that I can use for image generation that accurately describes the microstructure of a hierarchical composite.

Explicitly describe the geometry.}
\end{quote}

The comparison shown in Figure~\ref{fig:image_synth}(b)-(c) clearly highlights a more succinct prompt, which leads to a significantly more precise image generated by DALL-E 3~\cite{DALLECard}. Notably, the prompt produced by Zephyr-7B-$\beta$ fails to accurately capture the provided instruction and introduces additional design elements such as fibers and nanoparticles that were not mentioned in the original request focused strictly on the microstructure of a hierarchical composite. The X-LoRA prompt, being more concise and targeted, therefore yields a superior outcome.

\subsection{Generative solutions to protein analysis}

We now turn to specific generative and analytical tasks related to proteins. These capabilities arise within the X-LoRA model due to the protein mechanics adapter and are further strengthened by integrating the model's knowledge in biology, bio-inspired materials, reasoning, and logical deduction. This section focuses on analyzing protein-related tasks independently, evaluating the quantitative performance of this functionality.

Fig.~\ref{fig:protein_force_history_prediction} presents the outcomes from the protein task focused on predicting force-deformation behaviors based on sequences. Fig.~\ref{fig:protein_force_energy_prediction} illustrates the overall accuracy in forecasting unfolding force and unfolding energy from the sequence, comparing actual results with predictions (yielding $R_2=0.85$). In general, the model’s performance is excellent, and as will be discussed in the following section, it also performs exceptionally well for the inverse problem of design and cycle-consistent evaluation\cite{Lu2023GenerativePropertiesb}. However, unlike previous studies that relied on highly specialized GPT-type models for such tasks, our X-LoRA model integrates these specialized capabilities with a broad range of other deep specializations provided through various fine-tuned adaptations.  

\subsection{Protein design and analysis}

We now broaden the range of test cases to more comprehensive applications that involve multiple capability domains. In particular, the experiments described here demonstrate how the X-LoRA model can leverage a distinct set of skills and how agentic modeling can generate even more complex and thorough responses. We start by addressing a protein design task, where the model is tasked with designing a novel protein that exhibits a specific target force-deformation response. These force-deformation responses have been investigated in prior research~\cite{Ni2023ForceGen:Model} and represent measurements of force as a function of deformation when the protein is stretched from both ends (representing a classic protein unfolding experiment~\cite{Lu1998UnfoldingSimulation.}).

Fig.~\ref{fig:design_protein} illustrates the application of the generative protein task. We design a protein exhibiting a specified force-deformation response (based on training data and boundary conditions from the original molecular dynamics (MD) simulations of protein pulling, as detailed in~\cite{Ni2023ForceGen:Model}), then evaluate the predicted sequence to determine if it satisfies the desired target characteristics (Fig.~\ref{fig:design_protein}a). This cycle-consistent evaluation is crucial for verifying whether the inverse design capability aligns with the forward prediction.

In Fig.~\ref{fig:design_protein}b, the predicted sequence is displayed: \texttt{MSKSNTYRMLVLLEDDTKINKEDEKFLKGKPGKMHEFVDELILPFNVDELDELNTWFDKFDAEICIPYVEILKESGMK}. Additionally, Fig.~\ref{fig:design_protein}c presents a folded structural model of the protein (generated via AlphaFold 2~\cite{Jumper2021HighlyAlphaFoldb}). To investigate the relationship of the designed protein to existing sequences, Fig.~\ref{fig:design_protein}d shows its placement within a Basic Local Alignment Search Tool (BLAST) Tree~\cite{Altschul1990BasicTool}. Notably, the designed protein shares some affinity with hypothetical proteins previously studied but not yet structurally characterized or identified.

Employing agentic modeling with multiple X-LoRA models integrated into a multi-agent framework, we further examine the predicted sequence using our X-LoRA model. This approach leverages the enhanced ability of our model to comprehend protein mechanics and apply this knowledge in conjunction with other disciplines. Using the sequence predicted above as input, we request a more thorough analysis along with design recommendations for sequence variants. The initial task is:

\begin{quote}
\small\texttt{I have identified this amino acid sequence: MSKSNTYRMLVLLEDDTKINKEDEKFLKGKPGKMHEFVDEL ILPFNVDELDELNTWFDKFDAEICIPYVEILKESGMK.

How can I use it?

Please provide logical steps and specifically discuss the engineering aspects of the problem.}
\end{quote}

We employ multiple agentic interactions to construct the solution, utilizing adversarial strategies (see Materials and Methods for details) to encourage the model to explore more profound and intricate aspects of the concepts discussed. The summary of this dialogue, as generated by the X-LoRA model (full conversation available in Supplementary Material, SI Conv Example 1), is as follows:

\begin{quote}
\small\texttt{The discussion centers on the prospective uses of a particular amino acid sequence in protein engineering. This sequence includes both hydrophobic and hydrophilic residues, indicating it may form a stable conformation and engage with other molecules. To design a protein incorporating this sequence, computational approaches can be utilized to forecast its secondary and tertiary structures; however, challenges such as the complexity of protein folding and the high computational demands exist. Enhancements can be achieved by creating more efficient algorithms and integrating novel experimental data. Verifying computational models is essential to guarantee their precision and practicality. Strategies like cross-validation, regularization, early stopping, ensemble techniques, and transfer learning help avoid overfitting and underfitting. Moreover, interpreting and explaining models supports their application in real-world scenarios.}
\end{quote}

We also request the model to compile a summary of key points, producing:

\begin{quote}
\small\texttt{\begin{itemize}
    \item The amino acid sequence MSKSNTYRMLVLLEDDTKINKEDEKFLKGKPGKMHEFVDELI LPFNVDELDELNTWFDKFDAEICIPYVEILKESGMK holds promise in protein engineering owing to its combination of hydrophobic and hydrophilic residues. Computational tools can predict its secondary and tertiary structures, though they face challenges related to protein folding complexity and computational expense.
    \item Ensuring the accuracy and applicability of computational models is vital, using methods such as cross-validation, regularization, early stopping, ensemble learning, and transfer learning to mitigate overfitting and underfitting.
    \item Understanding and explaining model outputs further aids their practical use.
\end{itemize}
}
\end{quote}

Ultimately, the main conclusion is:

\begin{quote}
\small\texttt{The key insight is that the amino acid sequence MSKSNTYRMLVLLEDDTKINKEDEKFLKGKPGKMHEFVDELILPFNVDELDELNTWFDKFDAEICIPYVEILKESGMK holds promise for protein engineering applications due to its combination of hydrophobic and hydrophilic residues. To design a protein incorporating this sequence, computational techniques can be employed to predict its secondary and tertiary structures; however, challenges include the complexity of protein folding and the computational resources required. It is essential to validate computational models to ensure their accuracy and practical relevance. Approaches such as cross-validation, regularization, early stopping, ensemble methods, and transfer learning can be implemented to avoid overfitting and underfitting. Additionally, interpreting and explaining these models can aid in guiding practical applications.}
\end{quote}

We can further delve into detailed steps of sequence engineering. For this, we prompt the model to provide targeted design recommendations aimed at enhancing the protein's stability:

\begin{quote}
\small\texttt{How could I modify the following protein sequence to improve its stability: MSKSNTYRMLVLLEDDTKINKEDEKFLKGKPGKMHEFVDELILPFNVDELDELNTWFDKFDAEICIPYVEILKESGMK? }
\end{quote}

The summary of the dialogue is (full transcript available in Supplementary Material, SI Conv Example 2):

\begin{quote}
\small\texttt{The discussion centers on strategies to enhance protein sequence stability by incorporating hydrophobic amino acids and disulfide bonds. The biologist notes that hydrophobic residues generally promote protein stability, while disulfide bonds can also contribute to stability. The engineer inquires about specific hydrophobic amino acids and the optimal number of disulfide bonds to include. The biologist cites examples of proteins stabilized using these modifications, such as insulin, bovine pancreatic trypsin inhibitor, and green fluorescent protein. The conversation also addresses possible drawbacks of adding hydrophobic amino acids or disulfide bonds, including altered protein function, increased structural complexity, and the risk of aggregation. The biologist recommends mitigating these issues through careful amino acid selection, computational modeling, and experimental validation.}
\end{quote}

The summary of key points:

\begin{quote}
\small\texttt{\begin{itemize}
\item To enhance the stability of a protein sequence, hydrophobic amino acids and disulfide bonds may be incorporated. 

Hydrophobic amino acids generally contribute to increased protein stability, and disulfide bonds further strengthen this effect.

Proteins such as insulin, bovine pancreatic trypsin inhibitor, and green fluorescent protein serve as examples where these modifications have successfully improved stability.
\item Potential challenges or drawbacks of adding hydrophobic amino acids or disulfide bonds include possible alterations in protein function, greater structural complexity, and the risk of aggregation.
\item These challenges can be mitigated by careful amino acid selection, computational modeling, and experimental testing.
\end{itemize}
}
\end{quote}

In conclusion, the main point is:

\begin{quote}
\small\texttt{The most critical insight from the discussion is that to increase the stability of the protein sequence MSKSNTYRMLVLLEDDTKINKEDEKFLKGKPGKMHEFVDELILPFNVDELDELNTWFDKFDAEICIPYVEILKESGMK, one can introduce hydrophobic amino acids and disulfide bonds. This directly addresses the initial question.}
\end{quote}

From the perspective of protein engineering, these recommendations are logical, as they would indeed promote more stable protein designs. Specific examples and methods are proposed for researchers to further investigate.

In another scenario, we analyze three sequences. Initially, we request the model to evaluate the strength of each sequence, and then to reason about the outcomes: identify the strongest candidate and explain why that protein is likely the most robust. Additionally, we ask the model to describe the steps required to synthesize it in the laboratory. The three sequences (manually altered from one in the test set) are:

\begin{quote}
\small\texttt{
\begin{itemize}
    \item LNTWFDKFDAEICIPYVEILKESGMK \item AITWFDKFDAEICIPYVEALKIAAMK \item AAAAAAAAKFDAAEICIIAAAAAAAAKIAAMK
\end{itemize}
}
\end{quote}

This example evaluates whether the model can integrate multiple abilities and adapt dynamically as the conversation develops—that is, to flexibly switch between expertise such as protein mechanics, logical analysis, hypothesis generation, and so forth.

Figure~\ref{fig:conversation_evolve} provides a summary of the conversation, accompanied by an illustration of how the scalings evolve throughout the discussion (the analysis includes both the detailed deep layer-wise adapter scaling and a cumulative assessment to gauge the effective utilization of the different experts). The findings clearly indicate that the protein task is initially highly significant, progressively shifting to emphasize the bioinspired adapter and the biology adapter more strongly. By the conversation's conclusion, the biology adapter emerges as the dominant one.

To evaluate the predictions and assess the response quality, the three proteins involved in this task are shown in Figure~\ref{fig:protein_visuals} (proteins folded using Alpha Fold 2~\cite{Jumper2021HighlyAlphaFoldb}, via ColabFold \cite{Mirdita2022ColabFold:All}). It is observed that the most stable structure at the center exhibits the most well-organized secondary structure, consistent with the idea that it corresponds to the highest unfolding force, as forecasted by the X-LoRA model in Figure~\ref{fig:conversation_evolve}. The other two proteins are less structured; the first one is predominantly helical with some turns, while the last one is partially unstructured and displays a kink in its geometry. These structural characteristics account for the lowest predicted unfolding force as indicated by the model.

By examining the visual depictions of the folded conformations, we verify that the central, most stable structure exhibits the highest level of organized secondary structure, consistent with the idea that it corresponds to the greatest unfolding force as forecasted by the model. We observe that the other two structures display less order: the first is predominantly helical with some turns, while the last is partially disordered and contains a bend in its geometry. We propose that these characteristics probably account for the lowest unfolding force predicted by the model, aligning with the X-LoRA model's predictions.

In a second case study, we concentrate on the energetic characteristics of proteins, assessing the unfolding energy as the protein is stretched from its native state to a fully extended conformation~\cite{Ni2023ForceGen:Model}. We again analyze three sequences (the first two being the same as those previously considered, while the third is a longer construct combining the second sequence at the termini and the first sequence in the central region):

\begin{quote}
\small\texttt{\begin{itemize}
    \item LNTWFDKFDAEICIPYVEILKESGMK \item AITWFDKFDAEICIPYVEALKIAAMK \item AITWFDKFDAEICIPYVEALKIAAMKLNTWFDKFDAEICIPYVEILKESGMKAITWFDKFDAEICIPYVEALKIAAMK
\end{itemize}
}
\end{quote}

Initially, we instruct the model to calculate the unfolding energy for each sequence, and subsequently to analyze the findings: Given the amino acid sequences, we request an explanation for why one protein is likely the most stable. Finally, we inquire about the possible functions that the most stable protein could perform.

Figure~\ref{fig:MJB_20} presents a transcript of a dialogue that combines multiple specialized experts and integrates their knowledge, demonstrated here by analyzing the unfolding energy of three proteins. By the conclusion of the exchange, the CoT/reasoning adapter emerges as the most dominant, as the X-LoRA model addresses a query requiring reasoning over the results to forecast likely protein structural characteristics and possible functions. Specifically, the mechanistic inquiry into the source of protein stability highlights the formation of hydrophobic interactions among nonpolar amino acids and hydrogen bonds among polar amino acids. The model identifies several hydrophobic amino acids (I, F, W, Y) and hydrogen-bonding amino acids (D, E, K) that play a role in stabilizing the protein.

Additionally, the model predicts that cysteine (C) residues in the sequence will promote disulfide bond formation, further enhancing protein structural stability. Notably, the crucial predictions regarding structural features explaining the protein's high stability are accurately pinpointed, as validated by the structural analyses depicted in Fig.~\ref{fig:MJB_21}. This figure’s analysis confirms the presence of all these features and clarifies why the other two proteins, which are shorter segments with variable helical domain stability, exhibit much lower unfolding energies.

\subsection{Adversarial agentic modeling to connect distinct scholarly disciplines and knowledge yielding ontological knowledge graph generation}

Given that our X-LoRA model possesses extensive interdisciplinary knowledge, we leverage it to explore links among diverse concepts, knowledge bases, and fields of expertise. In one experiment, we pose two queries to the model, each similarly structured but differing in focus. We then generate triplets to construct an ontological knowledge graph that distills the responses into more structured representations~\cite{Giesa2012CategoryDesign,Spivak2011ReoccurringAnalogies}. This resulting graph offers a cohesive understanding of the insights produced and visually maps the relationships between concepts in a clear and mechanistic way.

We employ two adversarial agents to address the question. Agent 1 assumes the role of a "critical physicist," while Agent 2 takes on the role of a "philosopher." Agent 1 initiates the query and is responsible for probing and generating follow-up questions based on each response from Agent 2. As the dialogue progresses, the discussion delves deeper into the topic, revealing multiple facets. The question is deliberately chosen to span various disciplinary domains to evaluate the model’s capacity to synthesize ideas and concepts.

The query is:

\begin{quote}
\small\texttt{Consider an abstract representation of a classical musical composition as a graph of interacting notes, chords, melodies, rhythms and so on.

What happens if we would apply pressure, up to a point of failure, on this abstract concept of music?

Provide logical steps, mathematical reasoning, and discuss specifically the physics of the problem.}
\end{quote}

The entire conversation is available in the Supplementary Material (see, SI Conv Example 3). For conciseness, we present only a summary and key insights here, capturing the core elements:

\begin{quote}
\small\texttt{The discussion explores applying pressure to an abstract representation of a classical musical composition, treating it as a physical system with interacting notes, chords, melodies, and rhythms. This pressure is quantified by stress, and the failure point is identified when the stress reaches a critical threshold. To validate the mathematical model’s predictions, a physical model of the abstract musical composition may be constructed, allowing experimental results to be compared with theoretical expectations. The limitations of the proposed mathematical framework include oversimplified interaction representations, static modeling, uniform distribution assumptions, and a linear correlation between stress and pressure. Enhancing the model’s precision and applicability would require incorporating more detailed and realistic descriptions, dynamic behaviors, non-uniform distributions, and nonlinear terms.

[...]

The principal insight is that applying pressure to an abstract representation of a classical musical composition can be conceptualized as altering the force field among interacting notes, chords, melodies, and rhythms. As pressure increases, the system’s internal stress rises, culminating at a critical point where the system fails and collapses under its own load. This explanation addresses the original question by providing a physical perspective on how pressure influences the abstract notion of music and how such effects can be measured in terms of stress.}
\end{quote}

Although the generated text offers valuable perspectives, structuring the data into an ontological representation can enhance interpretability~\cite{Buehler2023MechGPTModalities}. For graph construction, we combine the entire content of both conversations to build a knowledge graph. An example of such a graph is shown in Figure~\ref{fig:graph}. The resulting graph is segmented into communities via the Girvan-Newman algorithm~\cite{Girvan2002CommunityNetworks}, identifying a total of 12 communities. The figure presents both the comprehensive graph (Figure~\ref{fig:graph}a) and detailed views with node labels (Figures~\ref{fig:graph}b-c).

\subsection{Creation of X-LoRA-Gemma Integrating Protein, Chemical, Bio-inspired, and Materials Mechanics Abilities}
\label{gemma-section}

To validate that the proposed methodology is effective across base models with different architectures, we trained an additional X-LoRA model, this time utilizing the Gemma-7B-it model~\cite{Google/gemma-7b-itFace}. This model incorporates a set of four LoRA adapters, defined as follows:

\begin{enumerate}
    \item Bioinspired materials \item Mechanics and materials \item Protein mechanics tasks (including generative sequence-to-property and inverse functionalities) \item Quantum-mechanics based molecular properties QM9 (including generative SMILES-to-property and inverse functionalities; see Table~\ref{tab:table_QM9} for definitions and a summary of all properties calculated in that dataset)~\cite{Ruddigkeit2012EnumerationGDB-17,Ramakrishnan2015ElectronicSpace}
\end{enumerate}

The X-LoRA-Gemma model is developed in a manner similar to the original X-LoRA model but aims to incorporate distinct capabilities, notably the ability to predict a set of 12 quantum mechanical properties and to design molecules satisfying those 12 quantum mechanical properties. This new adapter is trained jointly with the other adapters listed above, which rely on the same datasets as used in the prior experiment. Note that while the majority of the paper centers on the Zephyr-based X-LoRA model, for clarity in this section we differentiate the models by referring to the original X-LoRA as X-LoRA-Zephyr and the Gemma-based variant as X-LoRA-Gemma.

Figure~\ref{fig:protein_force_energy_prediction} presents the overall performance of X-LoRA, showcasing predictions of unfolding force and unfolding energy from the sequence ($R_2 = 0.87$), as well as results for predicting quantum mechanical properties of molecules using the QM9 dataset ($R_2 = 0.96$). We observe excellent outcomes for both tasks, with performance in the protein design task slightly surpassing that of the original X-LoRA model.

Next, we illustrate the capabilities of X-LoRA-Gemma through a straightforward multi-agent design experiment. We analyze a small organic molecule, \texttt{CCCC12CC1C(=O)N2} (selected from the QM9 test set). This molecule is a bicyclic structure containing a ketone (=O) and an amine (N) within the ring; see Figure~\ref{fig:QM9_design}(a) for a visualization. We pose the following query to the X-LoRA-Gemma model:

\begin{quote}
\small\texttt{Here is a molecule in SMILES code: CCCC12CC1C(=O)N2

What modifications would be necessary to make this molecule more suitable for polymer design.}
\end{quote}

The response is:

\begin{quote}
\small\texttt{To enhance the utility of this molecule for polymer synthesis, you might consider:
\begin{enumerate}
\item Adding a reactive functional group, like an amine or hydroxyl group, to facilitate further modification and crosslinking.
\item Extending the carbon backbone to be longer and more flexible, which would improve the molecule's capability to form polymers with varied structures and properties.
\item Attaching a side chain containing a functional group, such as a carboxyl group, to enable covalent bonding with other molecules and polymers.
\item Incorporating a cyclic moiety, for example, a lactam ring, to produce a polymer that is more rigid and stable.
\item Introducing diverse side chains to generate a polymer with a wide range of properties and functionalities.
\item Varying the number of carbon atoms in the backbone to allow the synthesis of polymers with differing molecular weights and characteristics.
\item Adding a crosslinking site to the molecule to produce a polymer that is more stable and less chemically reactive.
\item Implementing a repeating unit structure to create a polymer with defined molecular weight and attributes.
\item Combining several of these approaches to design a versatile polymer suitable for multiple applications.
\end{enumerate}}
\end{quote}

Hereafter, we concentrate on the first recommendation.

\begin{quote}
\small\texttt{How can I enhance the overall reactivity of a molecule? I intend to use this molecule for constructing a polymer material.

The molecule under consideration, represented in SMILES notation, is: CCCC12CC1C(=O)N2

I can adjust the following properties:

Dipole moment, mu Polarizability, alpha Highest Occupied Molecular Orbital energy, homo Lowest Unoccupied Molecular Orbital energy, lumo HOMO-LUMO gap: The energy difference between HOMO and LUMO, gap Electronic spatial extent Zero-point vibrational energy Heat capacity at constant volume Internal energy at 0 K Internal energy at 298.15 K Enthalpy at 298.15 K Free energy at 298.15 K Which are the three most important properties to alter, and should I increase or decrease each of them?}
\end{quote}

The reply from X-LoRA-Gemma is:

\begin{quote}
\small\texttt{\begin{itemize}
    \item  HOMO-LUMO gap (reduction) \item  Dipole moment (enhancement) \item  Polarizability (enhancement)
\end{itemize}
Altering these properties can boost the overall reactivity of the molecule, making it more appropriate for polymer material synthesis. }
\end{quote}

We employ both the generative and predictive functions of X-LoRA-Gemma to achieve these modifications in molecular properties.

For the optimization, we utilize a two-agent configuration composed of two instances of X-LoRA-Gemma. One agent performs the generative task \texttt{GenerateMolecularProperties<...>} to propose candidate molecules, while the second agent executes \texttt{CalculateMolecularProperties<...>} to evaluate the difference between the desired and predicted properties. This cycle continues until we hit the predefined iteration limit or achieve a mean-squared error below a set threshold.

Figure~\ref{fig:QM9_design} presents the outcomes of this procedure. Panel (a) illustrates the target design alongside its divergence from the original molecule's properties. Panel (b) displays the top 20 candidate molecules produced by the algorithm, with panel (c) showing the distribution of mean-squared errors, and panel (d) plotting the error across the generated designs. The best candidate, \texttt{CC(C)C1=CC(=O)NC=C1}, aligns closely with the target, as depicted in panel (e). Lastly, panel (f) offers an optimized 3D molecular structure, obtained via UFF~\cite{Rappe1992UFFSimulations} optimization.

To evaluate the molecule’s reactivity, we also calculated the Gasteiger charges~\cite{Gasteiger1980IterativeCharges} (visualized in Figure~\ref{fig:QM9_design}(f)). The molecule includes an aromatic ring bearing a ketone group (C=O), which is more electrophilic than typical carbon-carbon or carbon-hydrogen bonds, thus making it prone to nucleophilic attack. The presence of nitrogen within the aromatic ring together with the ketone forms a pyridone ring system. This ring exhibits higher reactivity than a simple aromatic ring because the heteroatom (nitrogen) and the ketone introduce reactive centers for chemical transformations such as nucleophilic addition at the carbonyl carbon or electrophilic substitution on the ring. The nitrogen atom, as part of the heteroaromatic ring, influences the aromaticity and chemical reactivity by modifying the electron density and its distribution within the ring. Consequently, this nitrogen-containing ring serves as a potential site for reactions, especially those involving electrophilic attacks, due to the electron-donating character of nitrogen.

The carbon atom within the ring, adjacent to both an oxygen and a nitrogen atom and carrying a partial positive charge (+0.25), indicates it is somewhat deficient in electron density. This occurs because oxygen, being more electronegative, attracts electron density toward itself, resulting in the carbon having less electron density. This effect is partially offset by the neighboring nitrogen atom, which is less electronegative than oxygen but more electronegative than carbon. Nonetheless, the partial positive charge signifies that the carbon atom is a probable site for nucleophilic attack, where a nucleophile (an electron-rich species) would be drawn to this electron-poor carbon.

The nitrogen atom exhibits a partial negative charge (-0.33), suggesting it is comparatively electron-rich relative to its typical state. This can be attributed to its lone pair of electrons and the influence of the aromatic ring’s electron system. In aromatic heterocycles, the nitrogen’s lone pair can engage in the delocalized $\pi$-electron system, though the precise extent of this participation depends on the ring’s structure and the electronegativity of neighboring atoms. The partial negative charge implies that the nitrogen is more nucleophilic, meaning it has an increased tendency to donate electrons, either by forming bonds with electrophiles or undergoing protonation reactions. Additional quantum mechanical studies should be conducted to explore this further.

The carbon carrying the partial positive charge serves as a reactive site for nucleophilic attacks. Other molecules could interact with this site via mechanisms where nucleophiles target electron-deficient carbons.

The hydrogen atom bonded to the nitrogen, which has a partial positive charge of 0.17, represents a favorable site for hydrogen bond formation with other molecules or atoms. Hydrogen bonds are a kind of dipole-dipole interaction occurring when a hydrogen atom, bonded to a highly electronegative atom (e.g., nitrogen, oxygen, or fluorine), interacts with another electronegative atom possessing a lone pair of electrons. This characteristic allows the molecule to potentially act as a hydrogen bond donor, which is an important factor influencing the molecule’s solubility in water and other polar solvents, as well as its behavior in biological contexts.

The ketone group adjacent to the nitrogen-containing aromatic ring imparts distinctive reactivity that can be leveraged in polymer synthesis~\cite{McQuarrieSimonPhysicalChemistry,CareySundbergAdvancedOrganicChemistry,Varghese2022BeyondApplications,Varghese2022BeyondApplications}. This compound is classified as a lactam, due to the cyclic amide formed by the ketone and nitrogen within the ring. Lactams are well-recognized in polymer chemistry, especially in the production of polyamides, a polymer category widely utilized in various material applications. The ketone functionality can participate in diverse chemical reactions valuable for polymer chemistry, such as nucleophilic addition. Such reactions can be harnessed to elongate the polymer chain or to introduce side chains that alter the polymer’s characteristics. The nitrogen atom within the aromatic ring can notably influence the polymer’s electronic attributes and can engage in hydrogen bonding, thereby affecting properties like melting point, solubility, and mechanical behavior. N-heteroaromatic compounds also exhibit the capacity to coordinate with metal ions, enabling applications in fields including catalysis or materials possessing electronic or photonic functionalities. This molecule may also undergo polymerization through its amide (lactam) group; polyamides derived from such reactions are known for their durability, flexibility, and resistance to abrasion and chemical exposure, finding uses across a broad spectrum from textiles (e.g., nylon) to automotive components~\cite{Varghese2022BeyondApplications,Varghese2022BeyondApplications}.

Alkyl substituents can enhance the hydrophobic nature of the molecule, so we expect that polymers synthesized from monomers containing such alkyl groups will exhibit increased hydrophobicity. This, in turn, influences their solubility in various solvents as well as their interactions with water. Such properties may be beneficial for uses that demand water resistance or particular solubility traits. Additionally, alkyl groups can impact the polymer's mechanical characteristics. For example, they may function as plasticizers within the polymer matrix, thereby improving the polymer's flexibility. This can result in materials that are more flexible or possess greater impact resistance, depending on the alkyl groups’ length and branching pattern.

\subsubsection{Additional benchmarking} We present further benchmarking results for the X-LoRA-Gemma model, displayed in Figure~\ref{fig:perf_MMLU}, utilizing a subset of the MMLU benchmark that concentrates on chemistry, physics, and biology (fields relatively similar to those on which our model was trained)~\cite{Hendrycks2020MeasuringUnderstanding}. Our findings show that both X-LoRA models described in this study outperform their respective base models. Specifically, the X-LoRA models surpass the base models’ performance (Zephyr model: 54.6\% for the base model versus 56.6\% for the X-LoRA model; Gemma model: 40.6\% for the base model compared to 52.4\% for the X-LoRA-Gemma model). The highest overall accuracy is achieved by the X-LoRA-Zephyr model, although the two models perform comparably well.

\section{Conclusion}

Multimodal LLMs have numerous existing and prospective applications in scientific fields, alongside a significant demand for highly specialized expertise tailored to specific purposes. Beyond general language processing tasks, the ability to perform computational or generative functions, as illustrated here with protein mechanics, is essential~\cite{Luu2023BioinspiredLLM:Materials,Buehler2023MechGPTModalities,Buehler2023MeLMProblemsb}. Built upon open-source language models, X-LoRA allows for dynamic blending of expertise and knowledge across various domains, as demonstrated through several examples, including Figure~\ref{fig:conversation_evolve}. This investigation revealed how the X-LoRA model autonomously and adaptively evolves its use of adapters, combining certain adapters in a sophisticated manner throughout the LoRA layers. This intricate blending of LoRA layers consistently appears in all cases analyzed here, such as in Figures~\ref{fig:QA_weights_sample_0} and \ref{fig:QA_weights}. This also presents an intriguing challenge to better comprehend the fundamental mechanisms from a physics viewpoint, considering an LLM as a large-scale graph-forming system that leverages complex graph structures. Consequently, the deep-layer mixing of adaptation experts, as conducted here, modifies these mechanisms to generate effective solutions.

The introduced algorithm extends the traditional inference approach of a single forward pass by incorporating two forward passes. This process trains the model to initially ‘reflect’ on the question and consider how to reconfigure itself before producing an answer. This approach embodies a straightforward implementation of ‘self-awareness,’ enabling the model to adjust its own architecture to optimally address a task. Therefore, despite having a relatively modest parameter count of 7 billion, the model can reason across a broad range of scientific disciplines (including biological materials, mathematics, physics, chemistry, logic, mechanics, and more), substantially improving its ability to produce innovative solutions. Additionally, it may inspire the design of alternative model architectures, such as strategies that generalize reconfigurability to portions or the entirety of non-adapted models. In such a framework, the scaling head would dynamically adjust segments of a trained model or integrate multiple models. As observed with X-LoRA, these strategies can serve as powerful tools to expand or unify the capabilities of distinct models.

We conducted multiple performance evaluations, including a comparison with significantly larger models (Figure~\ref{fig:perf_comp_bioinspiredexam}), where X-LoRA demonstrated clear superiority. The comprehensive comparison presented in Table~\ref{tab:table_qa} revealed that X-LoRA significantly outperforms the base model across a wide range of tasks and application areas. Additional comparisons involved quantitative evaluations of numerical prediction tasks, such as protein mechanics and quantum-mechanical molecular properties, with X-LoRA achieving high accuracy scores characterized by $R_2$ values between 0.85 and 0.96. By way of comparison, these results surpassed previously reported outcomes, such as those from the MeLM model~\cite{Buehler2023MeLMProblemsb} for predicting protein unfolding force, where an $R_2$ of 0.78 was recorded. The proficiency in accurately addressing multiple prediction challenges, inverse design tasks, and advanced reasoning abilities was highlighted in the molecular design example (Section~\ref{gemma-section}). This case also demonstrated that X-LoRA can be effectively implemented using different base model architectures (Gemma versus Zephyr/Mistral). Future research may extend these findings by developing X-LoRA models for larger base architectures.

A limitation of the proposed method is the necessity of two forward passes, which can increase computational overhead: the first pass computes hidden states that feed into the X-LoRA scaling head, and the second applies $\Lambda$ to the adapters. However, the initial pass utilizes the X-LoRA scaling head, thereby leveraging the fundamental capabilities of the LLM, as opposed to independently acquiring knowledge through training a larger scaling head. For practical deployment in model-serving frameworks, separate key-value caches must be maintained for both the scaling and forward passes. To enhance inference efficiency, we created \texttt{Mistral.rs}, an LLM serving platform written in Rust~\cite{matsakis2014rust} that integrates X-LoRA and incorporates multiple performance optimization techniques. Nevertheless, a key benefit of our approach is its applicability to any model without necessitating internal architectural modifications. We regard this as an advantage, especially given the extensive resources accessible within the Hugging Face ecosystem.

One benefit of the proposed method is that it offers a straightforward way to integrate with any existing LLM (for example, since the X-LoRA code has been made publicly available, it should be readily applicable to any model within the Hugging Face ecosystem). Another benefit lies in the scaling head, which leverages the inherent strengths of the underlying LLM and learns sophisticated techniques for optimally combining the adapters to produce certain responses. During the training of X-LoRA, it is ideal to present complex examples of question-answer pairs or conversations that enable the model to discover the most effective combinatorial strategy. We conducted several analyses related to protein design and mechanics. For example, Figure~\ref{fig:MJB_20} and Fig.~\ref{fig:MJB_21} show a stability analysis of multiple protein sequences, including comparisons, reasoning about the predicted outcomes, and a mechanistic explanation that offers a bottom-up chemical basis for the predicted behavior. As illustrated in Fig.~\ref{fig:MJB_21}, this was directly validated through a structural examination of the protein’s folded 3D conformation. 

Regarding future directions, further research could investigate a broader range of adapter experts, expanding upon the initial expertise represented in our LoRA collection. The X-LoRA model outperforms both the base model alone and the base model equipped with a single adapter, enabling it to answer complex questions by leveraging the specialized knowledge embedded in the adapter set. There also appear to be synergistic effects among the different adapters, as suggested by the intricate mixing patterns of scaling weights across layers. This phenomenon could be explored in greater depth and compared with approaches such as SLERP that integrate models using diverse strategies. While we trained the X-LoRA scaling head using a subset of the training data—consisting of a few hundred samples drawn from each original adapter training set—more targeted construction of a training dataset tailored for this purpose could be developed. For instance, we observed that presenting multiple-choice questions tends to activate reasoning-focused adapters, since these were trained on such types of data. Additional prompting techniques may be necessary to help the scaling head grasp specific domain knowledge relevant to the questions being asked, possibly through the use of specialized system messages or other methods. In general, creating suitable training datasets represents an area worthy of further investigation, including the development of particular methods to effectively induce layer-wise scaling combinations. Our experiments tracking these scaling factors over extended conversations demonstrate that the model can dynamically adjust scaling mechanisms to optimally handle distinct tasks, which is a promising capability.

We identified intriguing activation patterns of experts across the layers, with insightful discussions presented in the paper. This analysis can be further extended and might provide valuable understanding of the benefits of mixing different model components or sets of layers. Such investigations would be especially compelling when applying the X-LoRA method to domains beyond protein mechanics and design, as demonstrated here. We defer this exploration to future research.

Another significant direction is the deeper examination of agentic modeling, exemplified here through two adversarially interacting X-LoRA models. Subsequent studies could involve more than two models to enrich interactions and introduce additional functionalities, potentially driving models toward novel generative capacities. This approach could also facilitate the creation of methods that combine physics-based modeling or other validation techniques, employing code-writing/executing agents or agents leveraging function calling and other processing strategies~\cite{Ghafarollahi2024ProtAgents:Learning}.

\section{Materials and Methods}

The X-LoRA codes and sample implementations are accessible at \url{https://github.com/EricLBuehler/xlora}. The model weights and datasets associated with the X-LoRA presented here, along with further examples, can be found at \url{https://huggingface.co/lamm-mit/x-lora}. The X-LoRA-Gemma model is hosted at \url{https://huggingface.co/lamm-mit/x-lora-gemma-7b}. The \texttt{Mistral.rs} inference engine is available at \url{https://github.com/EricLBuehler/mistral.rs}.

\subsection{Mathematical formulation of X-LoRA} We define the $\mathbf{scalings}$ as a matrix comprising token-level and layer-level coefficients for the LoRA adapters. The X-LoRA $\mathbf{scaling}$ $\mathbf{head}$ receives hidden states from the base model and processes them through linear fully-connected layers to predict the scalings, followed by a softmax function applied along the last dimension to guarantee that the scalings across all experts sum to one. The scaling head in X-LoRA incorporates ReLU activation functions and dropout layers interleaved between linear fully-connected layers. The hidden layer size can be configured to either successfully reduce the dimensionality from the hidden state to the scaling dimension or to implement a widening and narrowing scheme akin to the feed-forward layer within the transformer block.

The scalings are represented as a matrix $\Lambda \in \mathbb{R}^{s \times l \times n}$, where $s$, $l$, and $n$ correspond to the sequence length, the number of layers, and the number of adapters, respectively. For each layer, the relevant scalings are extracted as a matrix $\lambda \in \mathbb{R}^{s \times n}$. These scalings $\lambda$ are then applied to each LoRA adapter by broadcasting and multiplying the corresponding X-LoRA scaling $\lambda_i \in \mathbb{R}^{1 \times 1 \times s}$ with the input to the decomposition matrices $A$ and $B$.

Given $W_0 \in \mathbb{R}^{d \times k}$, $B_i \in \mathbb{R}^{d \times r_i}$, and $A_i \in \mathbb{R}^{r_i \times d}$, where the rank $r_i \ll {\rm min} (d,k)$ is unique to each adapter $i$, the forward pass of an X-LoRA layer is defined as:

\begin{equation}
    h = W_0 x + \sum_{i=1}^{n} B_i A_i (x \cdot \lambda_i) \alpha_i
\end{equation}

Representations of the scalings are shown in the main paper, visualized using either bar plots or the \texttt{contourf} function from Matplotlib~\cite{MatplotlibDocumentation}.

\subsection{X-LoRA implementation algorithm and code} The implementation chosen here prioritizes simplicity and the flexibility to readily adapt the method to various models, particularly those available within the Hugging Face ecosystem. The X-LoRA algorithm functions by:

\begin{enumerate}
    \item Keeping the base model and adapters fixed, then fine-tuning performance.
    \item Combining LoRA adapters based on predictions generated by the X-LoRA scaling head.
\end{enumerate}

All code is developed using Python and PyTorch~\cite{Paszke2019PyTorchLibrary}.

\subsubsection{Dual forward pass strategy for self-aware inference} Since X-LoRA needs to compute hidden states to forecast the scalings, we employ a dual forward pass technique. This involves a first pass where the model assesses the task to decide how to reconfigure itself, followed by the actual inference step. This introduces a form of `self-awareness' that allows X-LoRA to allocate additional computation time for reflection rather than immediate response.

We refer to the initial pass as the scaling pass and the subsequent one as the forward pass. During the scaling pass, the scalings $\Lambda$ are set to zero. We also experimented with assigning scalings $\Lambda$ a value of ${n}^{-1}$ and observed comparable results, implying stable performance (this parameter can be configured in the X-LoRA package). To facilitate easy application of X-LoRA across different models, we introduce a forward pass hook. This technique is crucial for linking the forward and scaling passes and is the main feature enabling compatibility with all models.

\begin{enumerate}
    \item $\mathbf{Scaling}$ $\mathbf{pass}$: Performed by running the base model with adapters fixed to constant values. The hidden states produced are then utilized by the X-LoRA scaling head to compute $\Lambda$.
    \item $\mathbf{Forward}$ $\mathbf{pass}$: Use $\Lambda$ to blend the adapters during the model’s forward pass. The resulting logits are output as the prediction of the X-LoRA model.
\end{enumerate}

As a cautionary note, sharing the same key-value cache for both the scaling pass and forward pass can disrupt the correct inputs to the scaling head, since the cache would persist and accumulate across the two passes. Therefore, we disable the key-value cache during both training and inference and have implemented appropriate management of this in \texttt{Mistral.rs}.

\subsubsection{Freezing weights and selectively enabling model segments for training} The base model and adapters remain frozen (i.e., their parameters are not updated during training), making the X-LoRA scaling head the sole trainable component within an X-LoRA model. The provided code also allows the possibility to unfreeze all adapters alongside the X-LoRA scaling head. Additionally, one could choose to train the entire model simultaneously—comprising the X-LoRA scaling head, adapters, and the underlying base foundation model—potentially applying different learning rates to each part. Moreover, future developments of the techniques introduced here might facilitate the creation of scaling functions that operate across two or more foundation models, enabling a traceable realization of SLERP-like modeling. This avenue remains open for subsequent investigation.

\subsubsection{X-LoRA layers} An X-LoRA layer functions by scaling and then aggregating the outputs from multiple adapters. Each X-LoRA layer wraps around the original LoRA layer, thereby gaining access to the adapters specific to that layer.

To apply X-LoRA efficiently to any given model, the X-LoRA algorithm replaces the original LoRA layers with newly created X-LoRA layers and injects their forward computation logic into the LoRA layer. This approach utilizes features of Python’s object system to facilitate the layer swapping process.

In brief, during the forward pass, each X-LoRA layer performs the following operations (refer to Figure \ref{fig:general_arch}):

\begin{enumerate}
    \item $\mathbf{Scaling}$: Scale the input from each LoRA adapter by its corresponding scaling factor $\lambda_i$.
    \item $\mathbf{Forward}$: Execute the forward pass of the wrapped LoRA layer.
\end{enumerate}

\subsection{Efficient inference using \texttt{Mistral.rs} implemented in Rust} To enhance inference speed, we created \texttt{Mistral.rs}, a large language model serving platform implemented in Rust~\cite{matsakis2014rust} that incorporates X-LoRA along with multiple optimizations aimed at boosting performance. The following techniques, listed in no specific order, have been incorporated:

\begin{itemize}
    \item LoRA adapter weight stacking: When all matrices $A_i$ and $B_i$ from each LoRA adapter share identical shapes, these adapter weight matrices can be stacked along their first dimension.

Mathematically, starting with $W_0 \in \mathbb{R}^{d \times k}$, for adapter matrices $A_i \in \mathbb{R}^{r \times d}$ and $B_i \in \mathbb{R}^{d \times r}$ where $r_i \ll {\rm min} (d,k)$, we concatenate these matrices to form $A \in \mathbb{R}^{i \times r \times d}$ and $B \in \mathbb{R}^{i \times d \times r}$. The $\alpha_i$ parameters are also incorporated during the initialization of the stacked matrices.

To implement scalings, we rearrange the dimensions so that $\lambda \in \mathbb{R}^{i \times 1 \times 1}$. These scalings are then applied using broadcast multiplication.

\item Non-granular scalings To minimize the frequency of running the scaling pass, we provide an option for non-granular scalings. This approach caches the scalings after $n$ completed runs, since the KV cache guarantees that the sequence length remains one for the rest of the request. Employing this strategy significantly reduces latency by lowering the number of invocations of the base model.

\item Fused Compute Unified Device Architecture (CUDA) kernels To enhance the forward pass efficiency, we utilize fused CUDA kernels. These fused kernels allow complex operations, such as Rotary Embedding (\cite{su2021roformer}), which would typically execute sequentially, to be combined into a single kernel that runs in parallel on a GPU. Our kernels are based on the vLLM implementation (\cite{kwon2023efficient}).

\item Quantization The inference engine supports quantization, enabling the use of a quantized base model. Initial experiments indicate that quantization down to 8 bits performs well even for models originally trained with \texttt{bfloat16}, for example.

\end{itemize}

\subsection{Datasets}

Each adapter is trained on specialized datasets, summarized in Table~\ref{tab:table_loradata} with references to their sources and original publications.

\subsection{Training strategy} Training proceeds in phases. We employ the Zephyr-7B-$\beta$ model~\cite{HuggingFaceH4/zephyr-7b-betaFace}, which is built upon the Mistral-7B model\cite{Jiang2023Mistral7B}, and train a sequence of adapters using the datasets described previously. The Zephyr-7B-$\beta$ tokenizer is utilized.

The initial eight adapters are trained using question-answer pairs as provided by the datasets (see Table~\ref{tab:table_loradata}).

The protein mechanics tasks follow a similar training approach, but utilize specific forward and inverse instruction sets. The bidirectional instruction tasks include:

\begin{quote}
\tiny {\texttt{CalculateForce< G E E C D C G S P ....> [0.310]\\
CalculateEnergy< G E E C D C G S P ....> [0.220]\\ 
CalculateForceHistory< G E E C D C G S P ....> [0.004,0.034,0.125,0.142,0.159,0.102,0.079,0.073,0.131, ...]\\
GenerateForce<0.310>\,[ G E E C D C G S P ....]\\
GenerateEnergy<0.220>\,[ G E E C D C G S P ....]\\
GenerateForceHistory<0.004,0.034,0.125,0.142,0.159,0.102,0.079,0.073,0.131, ...>\,[ G E E C D C G S P ....] }}
\end{quote}

These correspond to three distinct tasks: determining the peak unfolding force of a protein, calculating the unfolding energy, and obtaining the force-deformation history. Each of these forward tasks is paired with a corresponding inverse task of the same kind, which produces a protein sequence as the output—resulting in a total of six tasks.

The rank of each LoRA adapter is set to $r=16$ with a scaling factor $\alpha=16$. Each adapter is trained targeting five specific modules: \texttt{v\_proj}, \texttt{k\_proj}, \texttt{q\_proj}, \texttt{gate\_proj}, and \texttt{o\_proj}, using a dropout rate of 0.05. Typically, LoRA adapters are trained using relatively low ranks, such as between 4 and 32, since it has been demonstrated that increasing the rank generally does not enhance performance but does increase computational costs. A smaller rank leads to a more compact model with fewer parameters~\cite{hu2021lora}. The selected values for $r$, $\alpha$, and other parameters reflect common settings from prior literature known to perform well. Additionally, previous work by the author has explored various parameter choices and determined that the values applied here represent a sensible configuration~\cite{Luu2023BioinspiredLLM:Materials}.

The X-LoRA scaling head is trained using approximately 2,000 samples, which are selected as a subset from the original dataset employed for training the adapters (we also augmented the training data with additional GPT-4 generated multiple-choice questions). We utilize a \texttt{paged\_adamw\_8bit} optimizer with gradient clipping set at 0.3, a learning rate of $2\times 10^{-4}$ with a warmup phase, and perform four gradient accumulation steps with a batch size of one. The X-LoRA model is trained for roughly 10,000 iterations. At the end of training, the loss is approximately 0.28.

Given that the underlying Mistral architecture has 32 layers and adapters are inserted into 5 of the transformer layers, there are a total of 160 LoRA layers, each replaced by an X-LoRA layer. The X-LoRA scaling head generates scaling factors for each of these 160 LoRA layers (with nine LoRA experts, this yields a total of $160 \times 9 = 1440$ $\lambda_i$ values calculated for every token). The scaling head consists of a single linear layer followed by ReLU activations and a dropout rate of $p=0.2$, comprising roughly 5 million parameters. The codebase supports flexible configurations, allowing the construction and exploration of more complex deep classifier heads in other contexts.

Both LoRA adapters and the X-LoRA scaling head are optimized using next-token prediction combined with cross-entropy loss (we estimate the probability assigned by the model to the correct next token in a sequence based on the preceding tokens; the training objective is to maximize this probability as indicated in the training data). We implement token shifting, where the input sequence is shifted by one token to create the target sequence the model aims to predict. For instance, given the input sequence ``Proteins are fundamental to," the corresponding target sequence for training becomes ``are fundamental to biology." Consequently, the model endeavors to predict the subsequent token in the sequence.

As depicted in Figure~\ref{fig:hierarchical_concept}, during the training of the X-LoRA scaling head, only the scaling head parameters are updated while all other model parameters remain fixed. The original Zephyr tokenizer and prompt template are employed.

\subsection{X-LoRA-Gemma model} The X-LoRA-Gemma model is constructed similarly to the previously described X-LoRA model but is built upon the Gemma-7B-it model~\cite{Google/gemma-7b-itFace}, utilizing four adapters (refer to the main text for further details). These four adapters are trained on datasets related to mechanics and materials, protein mechanics, bioinspired materials, along with an additional quantum-mechanics-based molecular properties dataset QM9 (see Table~\ref{tab:table_QM9})~\cite{Ruddigkeit2012EnumerationGDB-17,Ramakrishnan2015ElectronicSpace}.

Each LoRA adapter has a rank of $r=16$ and a scaling factor $\alpha=16$. The training targets for each adapter encompass seven modules: \texttt{q\_proj}, \texttt{o\_proj}, \texttt{k\_proj}, \texttt{v\_proj}, \texttt{gate\_proj}, \texttt{up\_proj}, and \texttt{down\_proj}, with a dropout rate of 0.05. With four adapters and seven target modules per adapter, considering the Gemma-7B base model contains 28 layers, the total number of $\lambda_i$ values computed per token amounts to $28 \times 28 = 784$.

The first two adapters (bioinspired materials and mechanics of materials) incorporated into this model are trained using question-answer pairs provided by the respective datasets (see Table~\ref{tab:table_loradata}).

Protein mechanics tasks follow the training methodology described previously, involving bidirectional forward and inverse instruction sets. Additionally, an adapter trained on the QM9 dataset includes these two tasks:

\begin{quote}
\tiny {\texttt{CalculateMolecularProperties<CC(C)N1CC1C=O> [0.098,0.358,0.581,0.395,0.309,0.330,0.570,0.649,0.519,0.519,0.519,0.518]\\
GenerateMolecularProperties<0.098,0.358,0.581,0.395,0.309,0.330,0.570,0.649,0.519,0.519,0.519,0.518> [CC(C)N1CC1C=O]\\ 
...
}}
\end{quote}

These encompass two main tasks: computing molecular properties (the forward task) and the inverse task of generating a molecule with desired properties. All 12 properties included in the QM9 dataset are either predicted or designed for, respectively (refer to Table~\ref{tab:table_QM9}).  
This model incorporates a total of four adapters.  

The X-LoRA scaling head for this model consists of three layers, with the intermediate layer sized at $3072 \times {FF}_{\rm fac} = 12,288$, where the feed-forward multiplier ${FF}_{\rm fac} = 4$. The architecture uses linear layers activated by ReLU functions, includes a dropout rate of $p=0.2$, and contains approximately 47 million parameters.  

The X-LoRA-Gemma scaling head was trained on roughly 17,000 samples. These samples were drawn as a subset from the training dataset used for adapter training. Notably, although the X-LoRA-Gemma model included only four adapters, we utilized samples from the full training set of the aforementioned model, supplemented with data from the QM9 dataset. This approach was taken to explore the potential for the X-LoRA model to acquire additional functionalities during the scaling head training phase. Given that the loss converged well to 0.58 by the end of training and exhibited excellent performance even on tasks the model was not specifically trained for (e.g., chemistry, as illustrated by the molecule design example presented in the paper), we consider that this method can generally produce strong results.

We employed a \texttt{paged\_adamw\_8bit} optimizer with a gradient clipping threshold of 0.3, a learning rate of $1\times 10^{-4}$ combined with warmup, and four steps of gradient accumulation, using a batch size of one. The X-LoRA-Gemma model was trained for approximately 62,000 steps.

The original Gemma tokenizer and prompt template were utilized.

\subsection{Adversarial agentic modeling} Our adversarial agentic approach involves creating two X-LoRA agents. One agent specializes in posing questions, initiating the initial inquiry and continuously generating further questions to uncover difficulties and flaws in the answers. The other agent handles answering these questions, providing ongoing responses. This approach is implemented using \{Guidance\}~\cite{Guidance-ai/guidance:Models.}.

Regarding the example that connects music and mechanics presented in the text, the characteristics of the two agents are defined as follows. The first agent serves as the physicist and questioner:

\begin{quote}
\small\texttt{You are a critical physicist engaged in a discussion from the standpoint of physics. Keep your responses concise and consistently challenge statements provocatively.

As a creative thinker, you incorporate ideas from other disciplines and push conventional boundaries.}
\end{quote}

The second agent is a philosopher (in some examples, this role is adapted to represent a biology expert):

\begin{quote}
\small\texttt{You are a philosopher with expertise in biology, chemistry, and mathematics. You participate in a discussion.

Keep your replies brief yet precise and imaginative. You generate excellent ideas and novel lines of thought, always maintaining logical rigor.}
\end{quote}

The questioner is explicitly directed to take into account the question and preceding answers, and to formulate new, probing questions. This is implemented through a prompt structured as follows:

\begin{quote}
\small\texttt{Examine the following question and answer.

\#\#\# Question: <question>

\#\#\# Response: <response>

\#\#\# Instruction: Provide a SINGLE follow-up question that critically examines the response.

DO NOT provide an answer or commentary yet.

The single question is: }
\end{quote}

After the dialogue has advanced for $N$ exchanges, the model is tasked with 1) summarizing the main insights covered, 2) enumerating the key insights as bullet points, and 3) determining the primary conclusion.

The unprocessed conversation text serves as the foundation for subsequent analysis via graph construction.

\subsection{Knowledge graph generation} We employ Zephyr-7B-$\beta$ to extract triplets from text, following the methodology outlined in~\cite{Rahulnyk/knowledge_graph:QnA} with added capabilities inspired by the Llama Index graph creation method~\cite{Liu_LlamaIndex_2022}. Graphs are rendered with NetworX\cite{Networkx/networkx:Python} and Pyvis~\cite{WestHealth/pyvis:Graphs.}. One key instruction provided, aligned with the approach in~\cite{Liu_LlamaIndex_2022}, is:

\begin{quote}
\small\texttt{Your objective is to extract the ontology of terms mentioned within the provided context. These terms should represent the principal concepts relevant to the context.

Present your output as a JSON list. Each element in the list includes a pair of terms and the relation connecting them, formatted as follows: [...] }
\end{quote}

We direct the construction of triplets comprising nodes and edges as follows, following the methodology described in \cite{Liu_LlamaIndex_2022}:

\begin{quote}
\small\texttt{\begin{itemize}
            \item "node\_1": "A concept from the extracted ontology"
            \item "node\_2": "A related concept from the extracted ontology"
            \item "edge": "A concise description of the relationship between node\_1 and node\_2"
        \end{itemize}
        }
\end{quote}

We also provide the model with an example of ontological analysis, such as:

\begin{quote}
\small\texttt{
       Context: ```Spider silk is a strong natural fiber used to catch prey in a web.``` }
\end{quote}

We then supply sample triplets, for instance:

\begin{quote}
\small\texttt{\begin{itemize}
            \item "node\_1": "spider silk"
            \item "node\_2": "fiber"
            \item "edge": "is"
        \end{itemize}
        }
\end{quote}

The compiled graph data is then organized and segmented into communities via the Girvan-Newman algorithm~\cite{Girvan2002CommunityNetworks}.

\subsection{Visualization of molecular structures}

PyMOL~\cite{Schrodinger2015The1.8} is utilized to visualize and analyze the predicted protein structures. Various scripts and built-in functions of this tool identify specific protein features such as secondary structure elements, hydrophobic and hydrophilic regions, disulfide bonds, and hydrogen bonds.

\section*{Data Availability Statement} The code and data underpinning this study’s results are freely accessible at \url{https://github.com/EricLBuehler/xlora} and \url{https://huggingface.co/lamm-mit/x-lora}. Additional data sources used herein are cited (refer to Table~\ref{tab:table_loradata} for details).

\end{document}